% Cursor Agent stream-drop card. ZH body. No private ports, nodes, or credentials.
% Heading then \par before every tabularx.
\documentclass[11pt,a4paper]{ctexart}
\usepackage[margin=1.55cm,top=1.6cm,bottom=1.75cm]{geometry}
\usepackage{xcolor}
\usepackage{tikz}
\usetikzlibrary{arrows.meta,positioning,fit,calc}
\usepackage{booktabs}
\usepackage{tabularx}
\usepackage{colortbl}
\usepackage{array}
\usepackage{hyperref}
\usepackage{enumitem}
\usepackage{fancyhdr}
\usepackage{microtype}
\usepackage{listings}

\hypersetup{
  colorlinks=true,
  urlcolor=linkb,
  pdftitle={Agent stopped retrying --- 中文圈土方法与官方阶梯}
}

\definecolor{ink}{HTML}{1A1A1A}
\definecolor{muted}{HTML}{5C5C5C}
\definecolor{accent}{HTML}{1F4E3D}
\definecolor{linkb}{HTML}{1F4E79}
\definecolor{paper}{HTML}{F7F4EE}
\definecolor{yes}{HTML}{E8F0EA}
\definecolor{no}{HTML}{F6E8E6}
\definecolor{mid}{HTML}{F7F0E4}
\definecolor{unk}{HTML}{E8E6E1}
\definecolor{codebg}{HTML}{EFEBE3}

\pagestyle{fancy}
\fancyhf{}
\renewcommand{\headrulewidth}{0.4pt}
\renewcommand{\footrulewidth}{0.3pt}
\fancyhead[L]{\small\color{muted}Europe/Berlin $\cdot$ 2026-09-06}
\fancyhead[R]{\small\color{accent}Cursor Agent 断流}
\fancyfoot[L]{\small\color{muted}tmp/cursor-agent-retry.pdf $\cdot$ cursor-agent-retry}
\fancyfoot[R]{\small\color{muted}\thepage}

\setlist[itemize]{leftmargin=1.1em,itemsep=0.06em,topsep=0.08em}
\setlist[enumerate]{leftmargin=1.2em,itemsep=0.06em,topsep=0.08em}
\setlength{\parskip}{0.14em}
\setlength{\parindent}{0pt}
\setcounter{secnumdepth}{0}

\lstset{
  basicstyle=\ttfamily\footnotesize,
  backgroundcolor=\color{codebg},
  frame=none,
  columns=fullflexible,
  keepspaces=true,
  breaklines=true,
  aboveskip=0.4em,
  belowskip=0.3em
}

\newcommand{\blocktable}[1]{%
\par
{\footnotesize
\noindent
\begin{tabularx}{\linewidth}{@{}
  >{\raggedright\arraybackslash}p{0.22\linewidth}
  >{\raggedright\arraybackslash}X
  >{\raggedright\arraybackslash}p{0.16\linewidth}@{}}
#1
\end{tabularx}}%
\par
}

\newcommand{\gradetable}[1]{%
\par
{\footnotesize
\noindent
\begin{tabularx}{\linewidth}{@{}
  >{\raggedright\arraybackslash}p{0.18\linewidth}
  >{\raggedright\arraybackslash}X
  >{\raggedright\arraybackslash}p{0.14\linewidth}@{}}
#1
\end{tabularx}}%
\par
}

\begin{document}
\pagecolor{paper}
\color{ink}

{\LARGE\bfseries\color{accent} Agent stopped retrying}\\[0.12em]
{\large 不是服务端掐你，是本机到 Cursor 的长连接在存盘前断了}

\vspace{0.25em}
{\color{muted}\small
官方帖 \href{https://forum.cursor.com/t/bug-agent-stopped-retrying-repeatedly-on-windows-standalone-desktop-app/169376/5}{forum 169376/5} $\cdot$
网络说明 \href{https://cursor.com/help/troubleshooting/network}{cursor.com/help/troubleshooting/network} $\cdot$
日文同页 \href{https://cursor.com/ja/help/troubleshooting/network}{JA} $\cdot$
CSDN / 掘金 / 博客园 / Easton 土方法对照。
Retry 从上次已保存状态接着跑，不是从头重做。}

\section{这句话在说什么}

Windows 独立客户端里，Agent 正在读文件、跑工具，流突然静音。客户端没法安全续跑同一段未保存工作，就停下来。员工看 Request ID：服务端还在流，是\textbf{你这头的连接变安静}。独立桌面版和 VS Code 扩展是两套设置，所以同一网络上扩展能跑、桌面版反复停，并不矛盾。

原帖作者已经试过：杀进程、HTTP/1.1、清 \texttt{workspaceStorage}、\texttt{.cursorignore}、关硬件加速。\textbf{对这条报错没用}。那些是别的病（空终端、索引卡死）。

\section{先走阶梯，再叠土方法}

改 Network 或 \texttt{settings.json} 之后：任务管理器里把所有 \texttt{Cursor.exe} 杀干净再开。Reload Window 不够。

\begin{enumerate}
  \item \textbf{Cursor Settings $\rightarrow$ Network $\rightarrow$ Run Diagnostics}，把结果留着。
  \item \textbf{HTTP Compatibility Mode $\rightarrow$ HTTP/1.1}（等同 \texttt{cursor.general.disableHttp2: true}）。仍断再试 \textbf{HTTP/1.0}（短请求，同类帖里换节点后恢复）。
  \item \textbf{手机热点对照}几分钟。热点稳、当前网不稳，就是当前路径。
  \item 把代理和 HTTP 模式对齐 VS Code 里的 Cursor 扩展。
  \item Clash / Mihomo / v2rayN：\textbf{规则模式}，把 \textbf{mixed/HTTP 端口}填进 \texttt{http.proxy}。不要填 SOCKS 口。能不走 TUN / 系统全局就先不走。
  \item \texttt{proxySupport} 默认 override、地址为空 $=$ \textbf{裸连}。填口，或改成 \texttt{on}/\texttt{fallback}。
  \item 诊断里握手失败：系统 DNS 改 \texttt{1.1.1.1} 或 \texttt{8.8.8.8}。
\end{enumerate}

\begin{center}
\begin{tikzpicture}[
  font=\small,
  box/.style={draw=accent,rounded corners=2pt,align=center,inner sep=5pt,fill=yes,text width=3.15cm},
  bad/.style={draw=accent,rounded corners=2pt,align=center,inner sep=5pt,fill=no,text width=3.15cm},
  midb/.style={draw=accent,rounded corners=2pt,align=center,inner sep=5pt,fill=mid,text width=3.15cm},
  arr/.style={-{Latex},accent,thick}
]
\node[box] (d) {Run Diagnostics};
\node[box,right=0.55cm of d] (h) {HTTP/1.1\\再试 1.0};
\node[midb,right=0.55cm of h] (n) {热点对照\\对齐扩展设置};
\node[box,below=0.55cm of d] (p) {规则模式\\填 HTTP 口};
\node[midb,below=0.55cm of h] (s) {override+空代理\\改成 on 或填口};
\node[bad,below=0.55cm of n] (x) {不要先清缓存\\不要 disableHttp1SSE};
\draw[arr] (d) -- (h);
\draw[arr] (h) -- (n);
\draw[arr] (d) -- (p);
\draw[arr] (h) -- (s);
\draw[arr] (n) -- (x);
\end{tikzpicture}
\end{center}

\section{中文圈土方法：能用 / 慎用 / 别用}

\gradetable{
\rowcolor{yes}
\textbf{做法} & \textbf{出处与要点} & \textbf{判定} \\
\addlinespace[0.12em]
HTTP/1.1 面板 &
CSDN 排坑、DeepSeek 社区「有魔法仍无模型」、掘金锁区补充。校园网 / 公司网 / VPN 最常见。CSDN 写「不用重启」——\textbf{官方要求整进程退出}。 &
先做 \\
\addlinespace[0.08em]
IDE 只给 Cursor 走代理 &
CSDN / 博客园 / 掘金：Clash \textbf{规则模式}，不需要全局、不需要虚拟网卡。\texttt{http.proxy} 填 mixed 口（常见 7890；Verge 常见 7897）。海外节点；香港也可能锁区。 &
先做 \\
\addlinespace[0.08em]
\texttt{proxySupport} &
Easton：默认 override $+$ 空 \texttt{http.proxy} $=$ 浏览器能翻、Cursor 裸连。改 \texttt{on}/\texttt{fallback}，或手填后再 override。 &
先做 \\
\addlinespace[0.08em]
换 DNS / 换节点 &
Easton：HTTP/2 失败有时是解析到坏的 Cloudflare 边。Clash 文：长连接被差线路重置。论坛 170410：换 VPN 节点 $+$ HTTP/1.0。 &
对照后做 \\
\addlinespace[0.08em]
Clash \texttt{http-version: "1.1"} &
Easton 兜底：代理链自己协商 HTTP/2 时，在 \texttt{config.yaml} 强制 1.1。 &
对照后做 \\
\addlinespace[0.08em]
PAC 把 \texttt{cursor.sh} 直连 &
少华：域名没墙 $\rightarrow$ 真实 IP $\rightarrow$ 锁区。规则里强制 \texttt{cursor.sh} 走代理。残留系统代理指向已死的 7897 也会 \texttt{ERR\_PROXY\_CONNECTION\_FAILED}。 &
登录/锁区时 \\
\rowcolor{mid}
\addlinespace[0.08em]
\texttt{proxyStrictSSL: false} &
CSDN 标配。公司 SSL 检查时才需要。会关掉证书校验。 &
慎用 \\
\rowcolor{mid}
\addlinespace[0.08em]
启动参数 &
Easton：旧版 \texttt{--disable-http2}；企业 \texttt{--proxy-auto-detect}。\texttt{--ignore-certificate-errors} 最后才用。Windows 可用 \texttt{Cursor.exe --no-proxy-server} 测系统代理残留。 &
慎用 \\
\rowcolor{no}
\addlinespace[0.08em]
\texttt{disableHttp1SSE: true} &
CSDN / 博客园当企业网「强化」。官方回退路径\textbf{就是} HTTP/1.1 SSE。先关 SSE 等于拆掉兼容层。 &
别先用 \\
\rowcolor{no}
\addlinespace[0.08em]
清缓存 / 关 GPU &
原帖已试，独立客户端仍复现。\texttt{workspaceStorage} 删的是聊天与工作区状态，治的是空终端，不是这条断流。 &
别先用 \\
}

\section{一份可粘的 settings.json}

端口只是例子。打开代理软件看 \textbf{mixed / HTTP} 口，不要填 SOCKS（v2rayN 的 10808 常被填错）。HTTP 与 HTTPS 都写。

\begin{lstlisting}
{
  "http.proxy": "http://127.0.0.1:7890",
  "https.proxy": "http://127.0.0.1:7890",
  "http.proxySupport": "override",
  "http.proxyStrictSSL": false,
  "cursor.general.disableHttp2": true
}
\end{lstlisting}

CLI：\texttt{\%USERPROFILE\%/.cursor/cli-config.json} 里加 \texttt{"network": \{"useHttp1ForAgent": true\}}，然后重启 \texttt{cursor-agent}。

\section{端口别混}

\blocktable{
\rowcolor{yes}
\textbf{软件} & \textbf{典型 mixed/HTTP} & \textbf{SOCKS} \\
\addlinespace[0.1em]
Clash / Mihomo & 7890 填这个 & 7891 不要填进 http.proxy \\
Clash Verge & 7897 Mixed Port & --- \\
v2rayN & 10809 HTTP & 10808 SOCKS \\
}

\section{防火墙要放的域名}

\texttt{*.cursor.sh}（含 \texttt{api2.cursor.sh}、\texttt{agent.api5.cursor.sh}）、\texttt{*.cursor-cdn.com}、\texttt{*.cursorapi.com}。企业代理还须：长连接不被强拆、SSE 不被缓冲。诊断若写 \texttt{HTTP/1.1 SSE responses are being buffered}，换节点或让网管关缓冲，再叠客户端开关也没用。

\section{日志与上报}

Windows 日志：\texttt{\%APPDATA\%/Cursor/logs/main.log}。搜 \texttt{ERR\_HTTP2}、\texttt{ECONNREFUSED}、证书失败、\texttt{ETIMEDOUT}。聊天菜单 Copy Request ID。诊断加 Request ID 才能让官方看路径，而不是猜版本。

\section{来源（中文优先）}

\begin{itemize}
  \item CSDN 排坑：\url{https://blog.csdn.net/weixin_41496173/article/details/149153599}
  \item CSDN 内置代理、关 TUN：\url{https://blog.csdn.net/weixin_66397563/article/details/155617027}
  \item 博客园同思路：\url{https://www.cnblogs.com/jzssuanfa/p/19503708}
  \item 掘金锁区 7897：\url{https://juejin.cn/post/7529020018618695730}
  \item Easton 中文排查：\url{https://eastondev.com/blog/zh/posts/dev/20260119-cursor-network-fix/}
  \item Easton 企业证书：\url{https://eastondev.com/blog/zh/posts/ai/20260529-cursor-proxy-config/}
  \item Clash 与长连接：\url{https://clashopen.com/zh-CN/blog/articles/cursor-clash-proxy-ai-ide-2026.html}
  \item 日文官方网络页：\url{https://cursor.com/ja/help/troubleshooting/network}
\end{itemize}

{\color{muted}\small 土方法能接上线，不保证安全。\texttt{proxyStrictSSL: false} 与忽略证书只在你知道公司在做 TLS 检查时用。本页不写私人节点、账号、PAC URL。}

\end{document}
